%==============================================================================
% example.tex — Complete usage example for the infufrgs class
%
% This file is hereby placed in the public domain.
%
% Compile sequence:
%   pdflatex infufrgs-example
%   bibtex   infufrgs-example
%   pdflatex infufrgs-example
%   pdflatex infufrgs-example
%==============================================================================

% Select program (ppgc, pgmicro, cic, ecp, ppgl) and document type
% (diss/mestrado, tese/doutorado, tc/dipl, ti, rp, espec, pep, prop-tese,
%  plano-doutorado).  Add `english' for English-language documents.
\documentclass[ppgc,diss,english]{infufrgs}

% If using pdfLaTeX, declare the input encoding (UTF-8 is recommended).
% Not needed with XeLaTeX or LuaLaTeX.
\usepackage[utf8]{inputenc}

% For figure inclusion
\usepackage{graphicx}

% Optionally load math packages
% \usepackage{amsmath}

%------------------------------------------------------------------------------
% Document metadata
%------------------------------------------------------------------------------
\title{Um Exemplo de Monografia do Instituto de Informática da UFRGS}
\translatedtitle{Using \LaTeX\ to Prepare Documents at INF/UFRGS}

% Author: \author{Surname}{First Name(s)}
\author{Flaumann}{Fritz Gutenberg}
% Multiple authors are supported:
%\author{Flaumann}{Frida Gutenberg}

% Advisor (title is optional): \advisor[Title]{Surname}{First Name(s)}
\advisor[Prof.~Dr.]{Lamport}{Leslie}
% Co-advisor is optional:
%\coadvisor[Prof.~Dr.]{Knuth}{Donald Ervin}

% Date of the defence.  Defaults to the current month and year if omitted.
%\date{maio}{2001}

% Location of the defence.  Defaults to Porto Alegre, RS.
%\location{Itaquaquecetuba}{SP}

%------------------------------------------------------------------------------
% Keywords — use Title Case; up to 10 keywords per language
%------------------------------------------------------------------------------
\keyword{Ciência da Computação}
\keyword{Formatação de documentos}
\keyword{LaTeX}

\translatedkeyword{Computer Science}
\translatedkeyword{Document formatting}
\translatedkeyword{LaTeX}

%------------------------------------------------------------------------------
% Nominata — override individual fields if the bundled defaults are stale.
% Example:
%   \renewcommand{\nominataReit}{Prof\textsuperscript{a}.~Nome Sobrenome}
%   \renewcommand{\nominataReitname}{Reitora}
%------------------------------------------------------------------------------

%==============================================================================
\begin{document}
%==============================================================================

% Title page + CIP (or nominata page)
\maketitle

%------------------------------------------------------------------------------
% Optional front matter
%------------------------------------------------------------------------------
\begin{dedicatoria}
  À minha família.
\end{dedicatoria}

\begin{agradecimentos}
  Agradeço a todos que contribuíram para a realização deste trabalho.
\end{agradecimentos}

%------------------------------------------------------------------------------
% Abstract in the primary language
%------------------------------------------------------------------------------
\begin{abstract}
  Este documento demonstra o uso da classe \textit{infufrgs} para a preparação
  de monografias, dissertações e teses no Instituto de Informática da UFRGS.
  A classe fornece formatação automática da folha de rosto, sumário, listas de
  figuras e tabelas, referências no padrão ABNT e outros elementos
  pré-textuais e pós-textuais.
\end{abstract}

%------------------------------------------------------------------------------
% Abstract in the second language (translatedabstract)
%------------------------------------------------------------------------------
\begin{translatedabstract}
  This document demonstrates the use of the \textit{infufrgs} class for
  preparing monographs, dissertations and theses at the Institute of
  Informatics of UFRGS.  The class provides automatic formatting of the
  title page, table of contents, lists of figures and tables, ABNT-style
  references, and other front- and back-matter elements.
\end{translatedabstract}

%------------------------------------------------------------------------------
% Lists
%------------------------------------------------------------------------------
\tableofcontents

\listoffigures

\listoftables

% Optional: list of abbreviations
\begin{listofabbrv}{UFRGS}
  \item[ABNT] Associação Brasileira de Normas Técnicas
  \item[INF]  Instituto de Informática
  \item[PPGC] Programa de Pós-Graduação em Computação
  \item[UFRGS] Universidade Federal do Rio Grande do Sul
\end{listofabbrv}

% Optional: list of symbols
% \begin{listofsymbols}{$\alpha$}
%   \item[$\alpha$] Alpha coefficient
% \end{listofsymbols}

%==============================================================================
% Text body
%==============================================================================

\chapter{Introdução}

Este capítulo introduz o tema do trabalho.  Citações bibliográficas seguem a
norma ABNT; por exemplo, a linguagem \LaTeX\ foi descrita por
\cite{lamport1994latex}.

Figuras devem ser referenciadas pelo número (Figura~\ref{fig:exemplo}).

\begin{figure}[htbp]
  \centering
  \fbox{\rule{0pt}{4cm}\rule{6cm}{0pt}}
  \caption{Exemplo de figura com legenda centralizada.}
  \label{fig:exemplo}
\end{figure}

\section{Organização do Texto}

O restante deste documento está organizado da seguinte forma.  O
Capítulo~\ref{cap:fundamentacao} apresenta a fundamentação teórica.
O Capítulo~\ref{cap:conclusao} apresenta as conclusões.

\chapter{Fundamentação Teórica}
\label{cap:fundamentacao}

\section{Formatação}

A classe \texttt{infufrgs} herda da classe \texttt{report} do \LaTeX.  As
margens, o espaçamento entre linhas, a numeração de páginas e os títulos de
capítulos e seções são configurados automaticamente de acordo com as normas
do Instituto.

\subsection{Tabelas}

Tabelas devem ter título acima do corpo e ser referenciadas no texto
(Tabela~\ref{tab:exemplo}).

\begin{table}[htbp]
  \centering
  \caption{Exemplo de tabela com legenda centralizada.}
  \label{tab:exemplo}
  \begin{tabular}{lcc}
    \hline
    \textbf{Coluna A} & \textbf{Coluna B} & \textbf{Coluna C} \\
    \hline
    Linha 1 & 1.0 & 2.0 \\
    Linha 2 & 3.0 & 4.0 \\
    \hline
  \end{tabular}
\end{table}

\subsubsection{Citações longas}

Citações com mais de três linhas devem ser destacadas em bloco, com recuo de
4\,cm à esquerda:

\begin{quote}
  Esta é uma citação longa de exemplo.  O ambiente \texttt{quote} da classe
  \texttt{infufrgs} aplica automaticamente o recuo, o tamanho de fonte e o
  espaçamento simples exigidos pela norma.
\end{quote}

\chapter{Conclusão}
\label{cap:conclusao}

Este documento demonstrou os principais recursos da classe \texttt{infufrgs}.

%==============================================================================
% Back matter
%==============================================================================

% Appendices start with \appendix (optional \annexname for annexes)
\appendix

\chapter{Exemplo de Apêndice}

Apêndices contêm material elaborado pelo próprio autor.

\annex

\chapter{Exemplo de Anexo}

Anexos contêm material de terceiros reproduzido para documentação.

%------------------------------------------------------------------------------
% References (ABNT author-date style)
%------------------------------------------------------------------------------
\bibliographystyle{abntex2-alf}
\bibliography{infufrgs-example}

\end{document}
